\documentclass[11pt,a4paper]{../_shared/altacv}

\geometry{left=1.5cm,right=1.5cm,top=1.cm,bottom=1.2cm,columnsep=1.5cm}

\usepackage[utf8]{inputenc}
\usepackage[T1]{fontenc}
\usepackage[french]{babel}
\usepackage{paracol}
\usepackage{xcolor}
\usepackage{hyperref}
\usepackage{fontawesome5}
\usepackage{setspace}

\definecolor{SlateGrey}{HTML}{2E2E2E}
\definecolor{LightGrey}{HTML}{666666}
\definecolor{DarkPastelRed}{HTML}{8AC0D9}
\definecolor{PastelRed}{HTML}{00B0DA}
\definecolor{GoldenEarth}{HTML}{B4AD0C}
\colorlet{name}{black}
\colorlet{tagline}{PastelRed}
\colorlet{heading}{DarkPastelRed}
\colorlet{headingrule}{GoldenEarth}
\colorlet{subheading}{PastelRed}
\colorlet{accent}{PastelRed}
\colorlet{emphasis}{SlateGrey}
\colorlet{body}{LightGrey}

\renewcommand{\familydefault}{\sfdefault}

\name{\Large BOILEAU Guillaume}
\personalinfo{
  \email{guillaume.boileaupro@gmail.com}
  \phone{+33 6 59 94 85 84}
  \location{Alpes Maritimes, France}
  \linkedin{guillaume-boileau-phd-891b81265}
  \researchgate{Guillaume-Boileau-2}
  \orcid{0000-0002-3576-6968}
}

\setlength{\parskip}{0.75\baselineskip}

\begin{document}
\makecvheader
\vspace{-0.6cm}
\cvsection{Lettre de motivation}
\vspace{-0.4cm}

{\color{accent}\textbf{Objet :}} \textit{Candidature au poste de Data Engineer senior}

{\color{body}

Madame, Monsieur,

Je vous adresse ma candidature au poste de Data Engineer senior au sein de REEL IT.

Votre annonce a immédiatement retenu mon attention, car elle correspond très directement à mon positionnement : concevoir, fiabiliser et industrialiser des flux de données complexes dans des environnements exigeants, avec une forte attente sur la performance, la robustesse et la qualité d’exécution. Je souhaite aujourd’hui mettre cette capacité au service d’une structure capable de porter des projets data ambitieux, avec une vraie culture du delivery et de la transformation.

Au cours de mon parcours, j’ai développé et coordonné des pipelines de traitement sur des environnements techniques complexes, avec un haut niveau d’exigence sur la qualité des résultats, la structuration des workflows, l’automatisation et la maintenabilité. J’ai l’habitude d’intervenir sur des chaînes de traitement impliquant ingestion, transformation, comparaison, validation et exploitation de données à grande échelle, en m’appuyant sur Python et sur une logique d’industrialisation forte.

Ce que je peux apporter à REEL IT, c’est d’abord une capacité à construire des pipelines fiables, lisibles et performants, mais aussi à sécuriser leur exécution dans la durée. Dans mes expériences précédentes, j’ai travaillé sur des systèmes où la résilience des traitements, la gestion des erreurs, la traçabilité, la reproductibilité et la qualité de la donnée n’étaient pas des options, mais des exigences de base. Cette rigueur m’a conduit à structurer les flux, documenter les traitements, améliorer les standards de développement et contribuer à des environnements de travail robustes et collaboratifs.

J’ai également une vraie culture de la performance et de l’optimisation. J’ai évolué dans des contextes où les volumes de données, la complexité des traitements et les contraintes de calcul imposaient une approche méthodique, orientée efficacité et passage à l’échelle. Cela m’a appris à concevoir des traitements robustes, à identifier les points de friction techniques, à améliorer les temps d’exécution et à maintenir un haut niveau de fiabilité sur l’ensemble de la chaîne de valeur de la donnée.

Mon expérience récente dans un environnement produit et services a renforcé cette approche très opérationnelle. J’y ai contribué au développement de services backend, à l’intégration de composants logiciels, au suivi de flux de données en conditions réelles et à la résolution d’incidents bloquants, avec une logique claire de continuité de service, de réactivité et de qualité d’exécution. Cette expérience complète utilement mon socle technique en apportant une dimension plus directement orientée exploitation, production et delivery.

Je sais également travailler dans des environnements transverses, au contact d’équipes techniques variées, avec la capacité de faire avancer les sujets de manière structurée et pragmatique. J’ai coordonné des activités techniques, encadré des contributeurs, organisé des priorités et assuré l’alignement entre besoins, implémentation et résultats. Cette capacité à conjuguer exécution technique, rigueur méthodologique et collaboration est au coeur de ma façon de travailler.

Je suis particulièrement intéressé par votre contexte autour d’une plateforme data à forte volumétrie, par les enjeux d’industrialisation de pipelines et par l’environnement Lakehouse que vous décrivez. Même lorsque les outils évoluent d’un contexte à l’autre, la logique reste la même : produire une donnée fiable, fraîche, exploitable, gouvernée, et garantir que les pipelines tiennent dans la durée. C’est précisément sur ce type d’enjeu que je souhaite apporter de la valeur.

Rejoindre REEL IT représenterait pour moi l’opportunité de contribuer à des projets data structurants, avec un haut niveau d’exigence technique et une vraie ambition de transformation. Je suis convaincu que mon parcours, à la croisée de l’ingénierie logicielle, des pipelines de données, de la coordination technique et de l’exécution rigoureuse, peut constituer un atout concret pour vos équipes.

Je serais heureux de pouvoir échanger avec vous afin de vous présenter plus en détail mon parcours et ma motivation.

Je vous prie d’agréer, Madame, Monsieur, l’expression de mes salutations distinguées.

\textbf{Guillaume Boileau}

}

\end{document}